\section*{Лекция 6. SVM}
В этой лекции мы покажем иной подход к задаче классификации, и постараемся ответить на насущный вопрос "--- а зачем нам нужен ещё какой-то подход, ведь казалось бы, итак уже есть логистическая регрессия?

Прежде чем перейти к сути метода опорных векторов, неплохо было бы вспомнить несколько базовых понятий, на которых он строится. Речь идёт об \emph{отступе} (margin) и о \emph{верхней оценке}, с помощью которой алгоритм принимает решение о принадлежности точки к тому или иному классу. 
Эти концепции являются фундаментальными для понимания того, почему SVM работает именно так, а не иначе, и почему он демонстрирует хорошую обобщающую способность на новых данных.

\subsection*{Разделяющая гиперплоскость}

Вообще, в классификации принято разделять типы выборок на \textit{линейно разделимую} и \textit{линейно неразделимую}. Из названия должно быть (надеюсь) понятно, что это означает, но давайте на всякий случай взглянем на следующее изображение:

\begin{figure}[H]
    \centering
    \includegraphics[width=0.75\linewidth]{img/lin_sep.png}
\end{figure}

Как можно увидеть, через данную выборку очевидным образом можно провести прямую, которая идеально разделит наши данные. Соответственно, линейно неразделимым случаем называется обратное.

Однако здесь появляется разумный вопрос "--- как нетрудно заметить, таких идеально разделяющих прямых можно построить бесконечно много, неужели они все будут считаться идеальными? И тут появляется следующий пример:

\begin{figure}[H]
    \centering
    \includegraphics[width=0.75\linewidth]{img/lin_sep_dop.png}
\end{figure}

Глазами видно, что оранжевая прямая хуже синей, но почему?

Дело в том, что оранжевая прямая вносит какие-то дополнительные объяснения данных. Например посмотрите на зелёную точку, которая интуитивно, очевидно, не должна принадлежать классу оранжевых звёздочек, но в силу такого построения она ему принаджелит. Задача сводится к построению не только разделяющей гиперплоскости, но и чтобы она была корректной. Тут и выручает SVM (Метод Опорных Векторов).

Здесь мы уже более явно работаем с \texttt{margin}, и явно стараемся максимизировать расстояние от точек до гиперплоскости (т.к. с точки зрения геометрии, такая прямая действительно будет лучшей для наших задач). Однако у обычного \texttt{margin} в этом смысле есть определённая проблема...

\subsection*{Геометрический отступ}

В классическом отступе, как мы ранее упомянали, не происходит нормирования расстояния, т.е. по факту это <<расстояние до гиперплоскости>> по факту может иметь нереалистичное значние. Мы его использовали, т.к. это расстояние оказалось наиболее интуитивно понятным и легко интегрируемым в дальнейшие модели. 

Для устранения зависимости от масштаба параметров вводится более устойчивая мера -- \emph{геометрический отступ}. Он определяется следующим образом:
\[
\gamma_i^{(\mathrm{geom})} = \frac{M}{\|w\|}.
\]

То есть по факту, теперь мы всё-таки берём более реальное расстояние. В контексте SVM нам это особенно пригодится.

\subsection*{Нормировка масштаба и постановка задачи}

Далее мы ставим перед собой следующую задачу:
\[
\min M = 1.
\]

А теперь тяжёлый для осознания трюк: мы полагаем, что минимальный допустимый \texttt{margin} будет единицей. Ранее мы упомянали формулу геометрического отступа, и теперь можем составить задачу его минимизации:
\[
\min \frac{ M  }{||w||}
\]

Смотря на условие выше, получится
\[
\min\frac{ M }{ ||w|| } = \frac{ 1 }{ ||w|| }
\]

Величину $\frac{ 2 }{ ||w|| }$ называют шириной разделяющей гиперплоскости.

Здесь остановимся немного подробнее: максимизацией геометрического отступа мы хотим добиться того, чтобы расстояние от точек до гиперплоскости было как можно больше. Это означает бОльшую уверенность классификатора в предсказании. 

Соответственно, здесь появляется первое отличие от логистической регрессии. Она делала упор на подсчёт вероятностей, а метод опорных векторов позволяет делать наиболее уверенные классификации.

Рассмотрим самый простой пример для наглядности. Пусть у нас всего два объекта в одномерном пространстве: \( x_1 = +1 \) с меткой \( y_1 = +1 \) и \( x_2 = -1 \) с меткой \( y_2 = -1 \).  
Возьмём простейший классификатор с параметрами \( w = 1 \), \( b = 0 \). Тогда функция оценки будет \( f(x) = x \), а разделяющая граница -- точка \( x = 0 \). Вычислим функциональные отступы:
\[
y_1 (w x_1 + b) = 1 \cdot (1 \cdot 1 + 0) = 1, \qquad y_2 (w x_2 + b) = (-1) \cdot (1 \cdot (-1) + 0) = 1.
\]
Оба объекта классифицированы правильно и находятся на одинаковом функциональном отступе, равном 1. Геометрический отступ (расстояние) для каждого из них составит \( 1 / |w| = 1 \). В данном случае любая другая разделяющая граница, расположенная между -1 и +1, также будет корректно классифицировать точки, но её геометрический отступ до ближайшей точки будет меньше. 
Например, если сместить границу в точку \( x = 0.5 \), то расстояние до ближайшей точки \( x_1 = 1 \) станет 0.5.
SVM же выберет именно ту границу, которая максимизирует этот минимальный отступ, то есть точку \( x = 0 \). 

Этот пример наглядно иллюстрирует, что задача SVM заключается не просто в нахождении любой разделяющей гиперплоскости, а в выборе такой, которая обеспечивает наибольший зазор между классами, проходя по наиболее безопасному, удалённому от данных участку.